\documentclass[11pt,a4paper]{article}
\usepackage[margin=1in]{geometry}
\usepackage{amsmath,amssymb,amsthm}
\usepackage{algorithm}
\usepackage{algpseudocode}
\usepackage{booktabs}
\usepackage{hyperref}

\theoremstyle{definition}
\newtheorem{definition}{\textbf{Definition}}[section]
\newtheorem{theorem}{\textbf{Theorem}}[section]
\newtheorem{lemma}[theorem]{\textbf{Lemma}}
\newtheorem{remark}{\textbf{Remark}}[section]

\title{\textbf{World Models and Recurrent State-Space Models (Dreamer)}}
\author{\textbf{Team Scaibu} \\ \small Email: \texttt{jaydeep.9664920749@gmail.com} \\ \small \textbf{Architecture and Core Engine Team}}
\date{\textbf{October 2026}}

\begin{document}
\maketitle

\section{Definitions and Cognitive Mental Models}

\subsection{Formal Mathematical Definition}
\textbf{Definition (Recurrent State-Space Models and Latent Imagination):}
Model-free reinforcement learning demands millions of physical environment interactions, which is prohibitively costly or dangerous in real-world robotics and web operations. \textbf{World Models and Dreamer} learn an explicit generative transition model of environment dynamics directly from past observations, enabling an actor-critic agent to train entirely inside the learned compact latent space (``in imagination'') without querying the physical world.

The core architecture is the \textbf{Recurrent State-Space Model (RSSM)}, which splits state representation into a deterministic recurrent component $h_t$ and a stochastic Markovian latent variable $z_t$:
\begin{enumerate}
    \item \textbf{Deterministic Recurrent Transition}:
    {\boldmath
    \begin{equation}
    h_t = f_\theta(h_{t-1}, z_{t-1}, a_{t-1})
    \end{equation}
    }
    \item \textbf{Stochastic Representation Posterior (Encoder)}:
    {\boldmath
    \begin{equation}
    q_\phi(z_t | h_t, x_t) \sim \mathcal{N}\left( \mu_\phi(h_t, x_t), \; \operatorname{diag}(\sigma_\phi^2(h_t, x_t)) \right)
    \end{equation}
    }
    \item \textbf{Stochastic Transition Prior (Dynamics Model)}:
    {\boldmath
    \begin{equation}
    p_\theta(z_t | h_t) \sim \mathcal{N}\left( \mu_\theta(h_t), \; \operatorname{diag}(\sigma_\theta^2(h_t)) \right)
    \end{equation}
    }
    \item \textbf{Observation and Reward Decoders}:
    {\boldmath
    \begin{equation}
    p_\theta(x_t | h_t, z_t) = \mathcal{N}(\hat{x}_t, I), \quad p_\theta(r_t | h_t, z_t) = \mathcal{N}(\hat{r}_t, 1)
    \end{equation}
    }
\end{enumerate}
The complete world model is trained by maximizing the Variational Evidence Lower Bound (ELBO) with KL balancing:
{\boldmath
\begin{equation}
\mathcal{L}_{\text{WM}} = \mathbb{E}_q \left[ \|x_t - \hat{x}_t\|^2 + (r_t - \hat{r}_t)^2 + \beta_{\text{scale}} D_{\text{KL}}^{\text{balanced}}(q_\phi(z_t) \,||\, p_\theta(z_t)) \right]
\end{equation}
}
Policy learning then generates multi-step rollouts purely in latent space $\tau = \{(h_\tau, z_\tau, a_\tau, \hat{r}_\tau)\}_{\tau=1}^H$, optimizing an actor-critic pair via generalized $\operatorname{TD}(\lambda)$ returns.

\subsection{Expanded Non-Technical Definition (Intuition for Anyone)}
\textbf{Intuitive Conceptual Definition:}
\begin{enumerate}
    \item \textbf{Dreaming Before Acting}: When a human learns to drive or play chess, they do not crash real cars 10,000 times. They close their eyes and simulate scenarios in their mind: ``If I turn sharply here, the car will spin.'' Dreamer gives neural networks that exact capability.
    \item \textbf{Compressing Reality}: High-definition camera images contain millions of raw pixels, most of which are irrelevant background. Dreamer's encoder compresses complex pixels into a tiny mathematical code (latent state) containing only the essential physics of what is happening.
    \item \textbf{A Sandbox in the Mind}: Inside this tiny mathematical sandbox, the agent can test millions of actions in seconds, predicting rewards and consequences instantly at zero real-world cost.
    \item \textbf{Sample Efficiency Miracle}: Because Dreamer can generate unlimited simulated practice rounds from a handful of real experiences, it learns complex robotic skills in hours instead of months.
\end{enumerate}

\subsection{Child-Friendly Mental Model (Physical Metaphor)}
Imagine building a toy race track in your bedroom before racing a real go-kart:
\begin{itemize}
    \item \textbf{The Brute Force Racer}: Gets in a real kart with blinders on, crashes into every wall 5,000 times, and slowly learns the track by getting bruised!
    \item \textbf{The Dreamer Racer}: Walks the real track once with a notepad. Then sits in their bedroom with a toy model track and practices 100,000 virtual laps in their imagination!
    \item \textbf{Stepping onto the Track}: When race day arrives, the Dreamer racer already knows every corner by heart and drives smoothly to victory on their very first real attempt.
\end{itemize}

\section{Core Mathematical Equations and Definitions}

\subsection{RSSM Variational KL Balancing}
{\boldmath
\begin{equation}
D_{\text{KL}}^{\text{bal}} = \alpha_{\text{bal}} D_{\text{KL}}\left( q_\phi(z_t | h_t, x_t) \,||\, \operatorname{sg}[p_\theta(z_t | h_t)] \right) + (1 - \alpha_{\text{bal}}) D_{\text{KL}}\left( \operatorname{sg}[q_\phi(z_t | h_t, x_t)] \,||\, p_\theta(z_t | h_t) \right)
\end{equation}
}
\begin{itemize}
    \item \textbf{Formal Definition}: Asymmetric KL divergence loss where stop-gradient ($\operatorname{sg}$) prevents the posterior from collapsing to an untrained prior.
    \item \textbf{What It Means}: Trains the prior dynamics to predict the posterior representations faster than the representations degrade toward trivial priors.
    \item \textbf{Why It Is Important}: Prevents catastrophic representation collapse in latent world models without requiring delicate heuristic warmup schedules.
\end{itemize}

\subsection{Analytic Diagonal Gaussian KL Divergence}
{\boldmath
\begin{equation}
D_{\text{KL}}(\mathcal{N}_q \,||\, \mathcal{N}_p) = \frac{1}{2}\sum_{i=1}^{d_z} \left( \ln\frac{\sigma_{p,i}^2}{\sigma_{q,i}^2} + \frac{\sigma_{q,i}^2 + (\mu_{q,i} - \mu_{p,i})^2}{\sigma_{p,i}^2} - 1 \right)
\end{equation}
}
\begin{itemize}
    \item \textbf{Formal Definition}: Exact closed-form Kullback-Leibler divergence between two diagonal multivariate normal distributions.
    \item \textbf{What It Means}: Measures the information divergence between the encoded perception posterior and predicted physics prior.
    \item \textbf{Why It Is Important}: Enables zero-variance analytical gradient computation during backpropagation through the world model.
\end{itemize}

\subsection{Generalized $\operatorname{TD}(\lambda)$ Latent Imagination Return}
{\boldmath
\begin{equation}
V_t^\lambda = \hat{r}_t + \gamma \left( (1 - \lambda) v_\psi(h_{t+1}, z_{t+1}) + \lambda V_{t+1}^\lambda \right)
\end{equation}
}
\begin{itemize}
    \item \textbf{Formal Definition}: Backward recursive target interpolating between 1-step bootstrapping and full Monte Carlo horizon returns in imagination.
    \item \textbf{What It Means}: Blends short-term critic estimates with longer imagined trajectory rewards to balance bias and variance.
    \item \textbf{Why It Is Important}: Serves as the primary learning target for training both actor and value networks inside Dreamer's imagined rollouts.
\end{itemize}

\section{Rigorous Mathematical Analysis Checklist}

\begin{itemize}
    \item \textbf{Notation}: $h_t$ (deterministic state), $z_t$ (stochastic latent), $x_t$ (observation), $r_t$ (reward), $\lambda$ (GAE factor), $\alpha_{\text{bal}}$ (KL balance parameter).
    \item \textbf{Epistemic Status}: Developed by Hafner et al. (Google Brain / DeepMind, 2019--2023 across DreamerV1, V2, V3).
    \item \textbf{Dependency Graph}: $(x_t, a_{t-1}) \to h_t \to (q(z_t), p(z_t)) \to (\hat{x}_t, \hat{r}_t) \to \mathcal{L}_{\text{WM}} \to \text{Imagination Rollout} \to V_t^\lambda \to \mathcal{L}_{\text{actor}}$.
    \item \textbf{Dimensions}: Observation space $\mathbb{R}^{d_X}$, latent state $\mathbb{R}^{d_Z}$, recurrent hidden dimension $d_H$, imagination horizon $H$.
    \item \textbf{Regimes}: Sample-efficient robotics, continuous simulation, visual Minecraft navigation, complex API control.
    \item \textbf{Invariants}: If $\mu_q \equiv \mu_p$ and $\sigma_q \equiv \sigma_p$, analytic KL divergence equals exactly $0.0$.
    \item \textbf{Isomorphisms}: RSSM is a non-linear Kalman filter with recurrent neural transition and variational autoencoder emission.
\end{itemize}

\section{Theoretical Theorems and Formal Proofs}

\begin{theorem}[Variational Lower Bound Guarantee of RSSM Dynamics]
The RSSM loss function $\mathcal{L}_{\text{WM}}$ optimizes a rigorous lower bound on the true observation log-marginal likelihood $\ln p_\theta(x_{1:T})$.
\end{theorem}

\begin{proof}
Let $z_{1:T}$ denote the latent sequence. By Jensen's inequality and importance weighting over variational posterior $q_\phi(z_{1:T} | x_{1:T}, a_{1:T})$:
{\boldmath
\begin{align}
\ln p_\theta(x_{1:T}) &= \ln \int p_\theta(x_{1:T}, z_{1:T}) \, dz_{1:T} \\
&= \ln \mathbb{E}_{q_\phi}\left[ \frac{p_\theta(x_{1:T}, z_{1:T})}{q_\phi(z_{1:T} | x_{1:T})} \right] \\
&\ge \mathbb{E}_{q_\phi}\left[ \ln \frac{p_\theta(x_{1:T}, z_{1:T})}{q_\phi(z_{1:T} | x_{1:T})} \right]
\end{align}
}
Factoring the joint prior and variational posterior according to the Markov structure of the RSSM:
{\boldmath
\begin{align}
p_\theta(x_{1:T}, z_{1:T}) &= \prod_{t=1}^T p_\theta(x_t | h_t, z_t) p_\theta(z_t | h_t) \\
q_\phi(z_{1:T} | x_{1:T}) &= \prod_{t=1}^T q_\phi(z_t | h_t, x_t)
\end{align}
}
Substituting these factorizations yields:
{\boldmath
\begin{equation}
\ln p_\theta(x_{1:T}) \ge \sum_{t=1}^T \mathbb{E}_{q_\phi}\left[ \ln p_\theta(x_t | h_t, z_t) \right] - \sum_{t=1}^T D_{\text{KL}}\left( q_\phi(z_t | h_t, x_t) \,||\, p_\theta(z_t | h_t) \right)
\end{equation}
}
Under Gaussian emission with unit variance, $\ln p_\theta(x_t | h_t, z_t) = -\frac{1}{2}\|x_t - \hat{x}_t\|^2 + \text{const}$.
Thus, minimizing the RSSM reconstruction error plus KL divergence strictly maximizes the variational lower bound on data log-likelihood.
\end{proof}

\begin{theorem}[Convergence of Backwards Recursive $\lambda$-Return]
For any bounded reward sequence $|\hat{r}_t| \le R_{\max}$ and bounded value estimates $|v_\psi| \le V_{\max}$, the recursive sequence $V_t^\lambda$ converges to a well-defined bounded target for all $\lambda \in [0, 1]$ and $\gamma \in [0, 1)$.
\end{theorem}

\begin{proof}
Let the recursion be defined for $t \in \{1, \dots, H-1\}$ with terminal boundary condition $V_H^\lambda = v_\psi(h_H, z_H)$:
{\boldmath
\begin{equation}
V_t^\lambda = \hat{r}_t + \gamma (1 - \lambda) v_\psi(h_{t+1}, z_{t+1}) + \gamma \lambda V_{t+1}^\lambda
\end{equation}
}
We prove boundedness by mathematical induction backwards from $H$.
At $t = H$, $|V_H^\lambda| \le V_{\max}$.
Assume for induction hypothesis that $|V_{t+1}^\lambda| \le \frac{R_{\max}}{1 - \gamma} + V_{\max}$.
Then for step $t$:
{\boldmath
\begin{align}
|V_t^\lambda| &\le |\hat{r}_t| + \gamma (1 - \lambda) |v_\psi| + \gamma \lambda |V_{t+1}^\lambda| \\
&\le R_{\max} + \gamma (1 - \lambda) V_{\max} + \gamma \lambda \left( \frac{R_{\max}}{1 - \gamma} + V_{\max} \right) \\
&= R_{\max}\left(1 + \frac{\gamma \lambda}{1 - \gamma}\right) + \gamma V_{\max} \le \frac{R_{\max}}{1 - \gamma} + V_{\max}
\end{align}
}
Because every term is bounded uniformly and the backwards linear recurrence is contractive with factor $\gamma \lambda < 1$, the return targets are stable and numerically non-divergent.
\end{proof}

\section{Foundational Architectural Questions and Epistemic Meta-Questions}

\subsection{Architectural Question 1: Latent vs Pixel Imagination}
\textbf{Question:} Why does Dreamer train its policy in latent space $z_t$ rather than decoding back to raw pixels $x_t$ at every step?\\
\textbf{Answer:} Decoding to pixels requires running expensive convolutional or transformer generative decoders at every step of rollout, which incurs massive latency. Rolling out transitions directly in compact latent space requires only small MLP evaluations, running over $100\times$ faster while discarding distracting visual background noise that does not affect rewards.

\textbf{Meta-Question:} How does the agent guard against exploiting flaws in the world model?\\
\textbf{Answer:} A learned model inevitably exhibits prediction drift over long horizons. Dreamer limits its imagination horizon to a moderate depth (typically $H \approx 15$ steps) and grounds the horizon tail using the learned critic $v_\psi$, preventing the actor from hallucinating infinite free rewards in regions of model inaccuracy.

\subsection{Architectural Question 2: KL Balancing vs Fixed Annealing}
\textbf{Question:} Why is KL balancing mathematically superior to simply scaling the KL loss with a fixed beta?\\
\textbf{Answer:} If the overall KL loss is scaled uniformly, the encoder posterior can change much faster than the dynamics prior, often collapsing representation variance to zero before the prior learns the physics. KL balancing applies a higher learning rate (weight $0.8$) to the prior, forcing it to chase the representation without destroying the encoder's descriptive capacity.

\textbf{Meta-Question:} Can Dreamer learn without any task rewards?\\
\textbf{Answer:} Yes. In zero-shot and unsupervised exploration regimes (such as Plan2Explore), the world model's prediction error or latent disagreement between an ensemble of transition models serves as an intrinsic exploration reward, allowing the agent to master environment physics before a downstream task is even defined.

\section{Algorithmic Implementation Specification}

\begin{algorithm}[H]
\caption{Dreamer RSSM Loss Evaluation and Latent Return Computation}
\begin{algorithmic}[1]
\Require Observations $\mathbf{x}$, reconstructions $\hat{\mathbf{x}}$, rewards $\mathbf{r}$, predictions $\hat{\mathbf{r}}$, posterior parameters $(\boldsymbol{\mu}_q, \mathbf{lv}_q)$, prior parameters $(\boldsymbol{\mu}_p, \mathbf{lv}_p)$, imagined rewards $\mathbf{r}_{\text{img}}$, values $\mathbf{v}_{\text{img}}$, weights $\beta_{\text{scale}}, \gamma, \lambda$
\Ensure Reconstruction loss $\mathcal{L}_{\text{rec}}$, reward loss $\mathcal{L}_{\text{rew}}$, KL loss $\mathcal{L}_{\text{KL}}$, total model loss $\mathcal{L}_{\text{WM}}$, returns $\mathbf{V}^\lambda$, actor loss $\mathcal{L}_{\text{act}}$
\State Initialize $rec\_err \gets 0.0, \quad rew\_err \gets 0.0, \quad kl\_accum \gets 0.0$
\For{$t = 0$ \textbf{to} $T-1$}
    \State $rec\_err \gets rec\_err + (\mathbf{x}[t] - \hat{\mathbf{x}}[t])^2$
    \State $rew\_err \gets rew\_err + (\mathbf{r}[t] - \hat{\mathbf{r}}[t])^2$
    \State $v_q \gets \exp(\mathbf{lv}_q[t]), \quad v_p \gets \exp(\mathbf{lv}_p[t])$
    \State $kl_t \gets 0.5 \cdot (\mathbf{lv}_p[t] - \mathbf{lv}_q[t] + (v_q + (\boldsymbol{\mu}_q[t] - \boldsymbol{\mu}_p[t])^2) / v_p - 1.0)$
    \State $kl\_accum \gets kl\_accum + \max(0.0, kl_t)$
\EndFor
\State $\mathcal{L}_{\text{rec}} \gets rec\_err / T, \quad \mathcal{L}_{\text{rew}} \gets rew\_err / T, \quad \mathcal{L}_{\text{KL}} \gets kl\_accum / T$
\State $\mathcal{L}_{\text{WM}} \gets \mathcal{L}_{\text{rec}} + \mathcal{L}_{\text{rew}} + \beta_{\text{scale}} \cdot \mathcal{L}_{\text{KL}}$
\State Allocate $\mathbf{V}^\lambda \in \mathbb{R}^H, \quad next\_v \gets \mathbf{v}_{\text{img}}[H]$
\For{$t = H-1$ \textbf{downto} $0$}
    \State $ret \gets \mathbf{r}_{\text{img}}[t] + \gamma \cdot ((1 - \lambda) \cdot \mathbf{v}_{\text{img}}[t+1] + \lambda \cdot next\_v)$
    \State $\mathbf{V}^\lambda[t] \gets ret, \quad next\_v \gets ret$
\EndFor
\State $\mathcal{L}_{\text{act}} \gets - \frac{1}{H}\sum_{t=0}^{H-1} \mathbf{V}^\lambda[t]$
\State \Return $(\mathcal{L}_{\text{rec}}, \mathcal{L}_{\text{rew}}, \mathcal{L}_{\text{KL}}, \mathcal{L}_{\text{WM}}, \mathbf{V}^\lambda, \mathcal{L}_{\text{act}})$
\end{algorithmic}
\end{algorithm}

\section{Big-$\Theta$ Complexity Analysis}

\begin{itemize}
    \item \textbf{World Model Loss Evaluation Time}: $\Theta(T)$ operations for sequence losses.
    \item \textbf{Analytic KL Divergence Time}: $\Theta(T)$ scalar evaluations of logarithm and exp.
    \item \textbf{Latent Imagination $\lambda$-Return Time}: $\Theta(H)$ backward recursive steps.
    \item \textbf{Space Complexity}: $\Theta(T + H)$ memory buffer for states and targets.
    \item \textbf{Parallel Reduction Depth}: $\mathcal{D} = \Theta(\log T)$ for world model error reductions.
\end{itemize}

\section{Single Canonical Repository Reference}

The complete deterministic scalar implementation, verification suite, and unit test benchmarks are published at:
\begin{center}
\url{https://github.com/Chief-Strategist-J/llm-obs-infra}
\end{center}

\end{document}
